\chapter{Resultados Esperados}
\label{cap:resultados}

Os resultados deste trabalho são de duas naturezas. Os resultados \textbf{técnicos} já foram obtidos: dizem respeito à construção e à validação de engenharia do artefato. Os resultados \textbf{pedagógicos} são esperados e dependem da avaliação com usuários, condicionada à aprovação no Comitê de Ética em Pesquisa. Espera-se que a Fase 1 do sistema permita coletar dados preliminares relevantes sobre a eficácia do uso da Realidade Virtual (RV) como ferramenta de treinamento em segurança na construção civil, no contexto brasileiro, com foco em trabalho em altura e uso correto de EPIs.

\section{Resultados Técnicos Já Obtidos}

A construção do artefato já produziu resultados de engenharia consolidados, detalhados no Capítulo~\ref{cap:validacao}: a arquitetura orientada a eventos foi implementada e validada de forma independente do motor de jogo, com suíte de testes executada fora do Unity e confinamento das dependências de SDK às camadas de borda. Esses resultados atestam a maturidade técnica do protótipo e viabilizam a etapa de avaliação com usuários, cujos resultados, descritos a seguir, são esperados.

\section{Eficácia do Treinamento}

Espera-se que o grupo experimental demonstre ganho de conhecimento significativo sobre procedimentos de segurança em trabalho em altura e uso correto de EPIs, evidenciado pela comparação entre pré e pós-testes e superior ao do grupo controle --- indicando que a abordagem imersiva contribui de forma diferenciada em relação ao vídeo instrucional.

\section{Motivação e Engajamento}

Espera-se que o grupo experimental relate motivação e engajamento mais elevados que o grupo controle, corroborando a literatura internacional sobre a superioridade da RV frente a métodos passivos de capacitação.

\section{Qualidade da Experiência Imersiva}

Espera-se que o sistema apresente usabilidade adequada no SUS --- acima de 68 pontos, a média histórica do instrumento --- e gere presença satisfatória no IPQ junto ao grupo experimental, indicando que é usável e imersivo o suficiente para uso como ferramenta de treinamento em contexto real.

\section{Contribuição Metodológica}

Independentemente das hipóteses principais, este estudo estabelece um protocolo de avaliação adaptado ao contexto brasileiro, com grupo controle e instrumentos padronizados aplicados a trabalhadores ou estudantes da construção civil sob as NRs nacionais. Os dados fornecerão os primeiros indicadores nacionais de tamanho de efeito para esse tipo de intervenção, viabilizando o dimensionamento amostral de estudos futuros com maior poder estatístico.

\section{Sobre o Teste de Retenção}

Caso o teste de retenção caiba no cronograma, espera-se que o grupo experimental apresente maior retenção do conhecimento após 30 a 45 dias que o grupo controle, indicando que a experiência imersiva favorece a consolidação do aprendizado a longo prazo.
